% !TeX root = ../main.tex
\section{Related Work}
\label{sec:related-work}

\subsection{Hardware-Aware Training and Robustness}

Hardware-aware training (HAT) mitigates analog non-idealities by injecting hardware perturbations into the forward path during optimization \citep{joshi2020accurate}. Although related in spirit to quantization-aware training \citep{choi2019pact} and to domain randomization in sim-to-real transfer \citep{tobin2017domain}, HAT for CIM faces an additional difficulty: device-to-device (D2D) mismatch is spatially structured and fixed for a given hardware instance. It is therefore different from independent thermal or sampling noise. Standard HAT usually keeps one random D2D mask fixed throughout training. Section 5.4 shows that this choice causes the model to fit a particular hardware instance rather than the deployment distribution. Ensemble HAT resamples the static D2D mask at each training epoch, improving zero-shot transfer to fresh hardware instances.

\subsection{Organic Neuromorphic and Optoelectronic Devices}
\label{subsec:organic-devices}

Organic synaptic and memristive devices have been widely explored as flexible, low-power candidates for neuromorphic hardware, with attractive features such as optical sensitivity, multilevel conductance tuning, and heterogeneous materials integration \citep{xu2025emerging,photonics2025organicreview}. Recent organic and hybrid optoelectronic studies have demonstrated multilevel memory, long retention tails, and synaptic behavior in several material stacks \citep{guo2024organic,jung2024organicfilaments,zeng2023organicmemristor}. More recent array demonstrations report energy-efficient organic device arrays, lithographically patterned near-infrared OPECT arrays, and large-scale optoelectronic synapse arrays for visual neural networks \citep{gebregiorgis2023organiccim,zhang2026opect,zhang2025mooptoelectronic,cui2025multimode}. At the same time, array-level studies are beginning to highlight active-matrix addressing, spatial optical response control, and crosstalk as practical system constraints \citep{visionarch2023crosstalk,amspa2024insensor,jang2023insensor}. Temperature-resilient organic synapses and high-temperature synaptic phototransistors similarly suggest that thermal stability is becoming relevant to deployment rather than remaining only a materials issue \citep{fuller2020tempresilient,guo2024hightemp}. Taken together, this literature provides plausible ranges for conductance windows, state counts, retention behavior, and array-level artifacts in an organic-CIM simulator.

At the same time, most of this literature remains device-centric. Network demonstrations are often limited to idealized mappings or small benchmarks, and variability is either absent or treated only qualitatively. The most relevant precursor for variability-aware organic simulation remains the Monte Carlo study of Alibart \emph{et al.}, which sampled device parameters for a simple organic-memristive perceptron \citep{alibart2016physical}. The present work follows that variability-aware perspective, but extends it to modern deep-learning software, CIFAR-scale tasks, and more expressive backbones.

\subsection{CIM Simulation Frameworks and Hybrid Mapping}
\label{subsec:cim-frameworks-mapping}

System-level CIM simulators such as DNN+NeuroSim established an effective template for connecting device assumptions, peripheral-circuit costs, and neural-network accuracy \citep{peng2020dnnneurosim}. MemTorch showed how memristive non-idealities can be embedded into PyTorch workflows for end-to-end evaluation \citep{lammie2022memtorch}, AIHWKIT exposed analog HAT and inference simulation in a practical mixed-signal software stack \citep{rasch2021aihwkit}, and CrossSim provides a GPU-accelerated crossbar-accuracy workflow with parasitic resistance, ADC, drift, and lookup-table-based training models \citep{crosssim2026}. However, these mature inorganic simulators cannot natively model inverse-gamma optoelectronic photoresponse or organic-specific double-exponential retention without massive structural hacking. Thus, our behavioral profile-driven approach serves as a necessary complement—specifically tailored for organic optoelectronics—rather than a direct competitor to these inorganic toolkits.

In the CIM context, dense linear operators are natural crossbar targets, whereas depthwise and highly irregular operators often suffer from poor utilization \citep{wu2023bwq,wang2024epim,ge2024allspark}. Recent heterogeneous CIM accelerators for Vision Transformers confirm this pattern: linear projections (QKV, output, feed-forward) are mapped to analog arrays, whereas attention scores, softmax normalization, and dynamic interactions usually remain digital because of their precision requirements and irregular access patterns \citep{kim2025hemlet,lin2024hardsea,bettayeb2024memristorattention}. This separation underlies our mapping policy. Query-key-value projections, output projections, feed-forward layers, and patch-embedding convolutions are treated as analog candidates, while depthwise convolutions, normalization, and dynamic attention products remain digital. We adopt this hybrid logic, extending prior CIM simulators to organic-device profiles and Tiny-ViT operators.

Unlike prior ViT-on-PIM studies, we explicitly model photoresponse, profile substitution, and joint variability, retention, ADC resolution, and front-end distortion within one workflow. This focus is also consistent with recent low-bit ViT quantization studies, which repeatedly identify attention logits, softmax-adjacent computations, and activation scaling as especially sensitive under aggressive precision reduction \citep{liu2021ptqvit,li2022qvit,lin2023vitptq}. Post-training quantization studies likewise note that attention remains difficult below 6 bits unless specialized approximations are introduced \citep{lin2023vitptq}. By extending established simulator practice to organic optoelectronic domains, we aim to provide a deployment-facing framework that links reported device characteristics to CIFAR-scale transformer benchmarks.
